\documentclass[11pt]{article}

\usepackage[utf8]{inputenc}
\usepackage[T1]{fontenc}
\usepackage[margin=1in]{geometry}
\usepackage{hyperref}
\usepackage{url}
\usepackage{booktabs}
\usepackage{amsfonts}
\usepackage{nicefrac}
\usepackage{microtype}
\usepackage{amsmath}
\usepackage{amssymb}
\usepackage{graphicx}
\usepackage{xcolor}
\usepackage{algorithm}
\usepackage{algorithmic}
\usepackage{times}
\usepackage{titlesec}

\hypersetup{
    colorlinks=true,
    linkcolor=blue,
    filecolor=blue,
    urlcolor=blue,
    citecolor=blue,
}

\titleformat{\section}{\large\bfseries}{\thesection}{1em}{}
\titleformat{\subsection}{\normalsize\bfseries}{\thesubsection}{1em}{}

\title{\textbf{G3P: Gradient-Guided Privacy-Preserving Pruning via Train-Test Gradient Saliency}}

\author{%
  Anonymous Author(s) \\
  Affiliation \\
  \texttt{email@domain.com}
}

\date{}

\begin{document}

\maketitle

\begin{abstract}
\noindent Neural networks face significant privacy vulnerabilities from membership inference attacks (MIAs) that determine whether specific samples were used during training. While model pruning has shown promise as a privacy defense, prior work relies on incidental effects rather than explicitly targeting privacy-leaking neurons. We propose G3P (Gradient-Guided Privacy-Preserving Pruning), a method that identifies and removes neurons contributing to membership inference vulnerability through a novel privacy saliency metric based on train-test gradient magnitude differences. Our key insight is that neurons exhibiting significantly different gradient behaviors between training and test data are the primary contributors to privacy leakage. G3P combines this privacy saliency with task importance to guide structured pruning. Experiments on CIFAR-10 with ResNet-18 demonstrate that G3P reduces MIA AUC from 0.582 (magnitude pruning) to 0.527 at 30\% sparsity---a 9.5\% reduction in attack success. Notably, G3P outperforms a hybrid baseline combining magnitude pruning with KL-divergence regularization (MIA AUC 0.537), demonstrating that gradient-based neuron selection provides privacy benefits beyond output regularization alone. These results establish train-test gradient differences as a principled criterion for privacy-preserving pruning.
\end{abstract}

\section{Introduction}

Deep neural networks have achieved remarkable success across numerous applications, but they face significant privacy vulnerabilities. Membership Inference Attacks (MIAs) pose a particularly concerning threat, enabling adversaries to determine whether a specific individual's data was part of a model's training set \citep{shokri2017membership}. The ability to infer membership can reveal sensitive information, such as medical diagnoses from healthcare datasets or personal preferences from behavioral data.

Standard privacy-preserving techniques like Differential Privacy (DP) provide strong theoretical guarantees but often suffer from severe utility degradation. DP-SGD \citep{abadi2016deep}, the most widely used DP training method, typically incurs 5--15\% accuracy drops for meaningful privacy budgets ($\varepsilon \leq 4$), limiting its practical deployment.

Model compression through pruning offers an alternative path to privacy preservation. \citet{wang2021against} demonstrated that magnitude-based pruning can reduce MIA attack accuracy by 10--13.6\%, suggesting pruning's potential as a privacy defense. However, \citet{yuan2022membership} showed that pruning can actually \emph{increase} memorization and privacy risks in certain scenarios. This inconsistency raises a critical question: \textbf{Can we design pruning strategies that consistently improve privacy by explicitly targeting privacy-leaking neurons?}

Our key insight is that not all neurons contribute equally to privacy leakage. We hypothesize that neurons exhibiting significantly different gradient magnitudes between training and test data are the primary contributors to membership inference vulnerability. The intuition is that these neurons ``memorize'' training-specific patterns, producing distinct gradient signals for member versus non-member data---precisely the signal exploited by MIAs.

Based on this insight, we propose \textbf{G3P} (Gradient-Guided Privacy-Preserving Pruning), a method that:
\begin{enumerate}
    \item \textbf{Quantifies neuron-level privacy contribution} using the absolute difference in gradient magnitudes between training and test distributions;
    \item \textbf{Combines task importance and privacy risk} into a unified saliency score;
    \item \textbf{Progressively prunes privacy-vulnerable neurons} while preserving task-critical ones;
    \item \textbf{Fine-tunes with privacy-aware regularization} to further reduce leakage.
\end{enumerate}

\textbf{Our contributions:}
\begin{itemize}
    \item We propose a novel \emph{privacy saliency metric} based on train-test gradient magnitude differences, providing a principled way to identify privacy-leaking neurons.
    \item We demonstrate that gradient-based neuron selection provides privacy benefits beyond combining magnitude pruning with KL-divergence regularization, addressing a key gap in prior work.
    \item We conduct extensive experiments showing that G3P reduces MIA AUC by 9.5\% compared to magnitude pruning while maintaining competitive accuracy.
\end{itemize}

\section{Related Work}

\subsection{Membership Inference Attacks}

\citet{shokri2017membership} introduced the first membership inference attacks using shadow models to train attack classifiers. \citet{salem2019mlleaks} proposed more efficient attacks requiring fewer assumptions. The current state-of-the-art is LiRA (Likelihood Ratio Attack) by \citet{carlini2022membership}, which evaluates true-positive rates at low false-positive rates ($\leq$0.1\%) using likelihood ratios derived from multiple shadow models. Recent work by \citet{zhang2024spurious} identified that certain subpopulations exhibit disproportionately higher privacy vulnerabilities---up to 100$\times$ more vulnerable than others---revealing privacy disparities that aggregate metrics hide.

\subsection{Privacy Defenses}

Differential Privacy \citep{abadi2016deep} provides formal privacy guarantees through gradient clipping and noise injection but suffers from significant utility degradation. Regularization-based defenses like RelaxLoss \citep{chen2022relaxloss} align training and test loss distributions without formal guarantees.

\subsection{Pruning and Privacy}

\citet{wang2021against} first demonstrated that magnitude-based pruning can reduce MIA vulnerability as a side effect of compression, providing theoretical analysis connecting pruning's regularization effect to privacy improvement. However, their approach does not explicitly target privacy-leaking neurons.

\citet{yuan2022membership} showed that pruning can increase privacy risks in some scenarios and proposed KL-divergence regularization as a defense. Critically, their approach applies KL regularization uniformly to all neurons during fine-tuning without modifying the pruning criterion itself.

PriPrune \citep{chu2024priprune} targets federated learning settings with pseudo-pruning masks to protect against gradient inversion attacks, differing from our centralized setting focus.

\textbf{Key Gap:} While prior work has shown that pruning \emph{can} improve privacy, no existing method explicitly identifies and removes privacy-leaking neurons through \textbf{privacy-specific gradient-based saliency metrics}. Our contribution is a principled criterion for selecting which neurons to prune, not merely the general idea of using pruning for privacy.

\section{Method}

\subsection{Problem Formulation}

Given a trained neural network $f_\theta$ with parameters $\theta$, our goal is to produce a pruned model $f_{\theta'}$ with reduced parameters that:
\begin{enumerate}
    \item Maintains task performance: $\mathcal{L}(f_{\theta'}) \approx \mathcal{L}(f_\theta)$
    \item Reduces membership inference vulnerability: $\text{MIA\_AUC}(f_{\theta'}) < \text{MIA\_AUC}(f_\theta)$
\end{enumerate}

We focus on structured channel pruning, which removes entire output channels from convolutional layers, enabling practical speedups on standard hardware.

\subsection{Privacy Saliency Metric}

Our core innovation is defining privacy saliency based on the divergence between training and test gradient behaviors. For each neuron (channel) $i$, we define:

\begin{equation}
S_{\text{privacy}}(i) = \left| \mathbb{E}_{(x,y) \sim \mathcal{D}_{\text{train}}} \left[ \left| \frac{\partial \mathcal{L}}{\partial h_i} \right| \right] - \mathbb{E}_{(x,y) \sim \mathcal{D}_{\text{test}}} \left[ \left| \frac{\partial \mathcal{L}}{\partial h_i} \right| \right] \right|
\end{equation}

where $h_i$ is the activation of neuron $i$ and $\mathcal{L}$ is the loss function. The intuition is that neurons with significantly different gradient magnitudes between training and test data contribute to the membership signal exploited by MIAs.

\subsection{Task Saliency Metric}

Following \citet{molchanov2019importance}, we estimate task importance using first-order Taylor expansion:

\begin{equation}
S_{\text{task}}(i) = \left| \mathbb{E}_{(x,y) \sim \mathcal{D}_{\text{train}}} \left[ \frac{\partial \mathcal{L}}{\partial h_i} \cdot h_i \right] \right|
\end{equation}

This approximates the change in loss if neuron $i$ were removed.

\subsection{Combined Saliency Score}

We combine task and privacy saliency into a unified score:

\begin{equation}
S_{\text{combined}}(i) = \alpha \cdot S_{\text{task}}(i) - \beta \cdot S_{\text{privacy}}(i)
\end{equation}

where $\alpha, \beta \geq 0$ are hyperparameters controlling the privacy-utility trade-off. Neurons with the lowest $S_{\text{combined}}$ are pruned first. Setting $\beta = 0$ recovers standard task-oriented pruning, while $\beta > 0$ explicitly prioritizes removing privacy-leaking neurons.

\subsection{Progressive Privacy-Aware Pruning}

G3P operates through an iterative process:

\begin{algorithm}
\caption{G3P: Gradient-Guided Privacy-Preserving Pruning}
\begin{algorithmic}[1]
\STATE \textbf{Input:} Trained model $f_\theta$, target sparsity $s$, hyperparameters $\alpha, \beta$, iterations $T$
\STATE Initialize: $\theta^{(0)} \leftarrow \theta$
\FOR{$t = 1$ to $T$}
    \STATE Compute $S_{\text{privacy}}(i)$ using Eq.~1 on train and validation sets
    \STATE Compute $S_{\text{task}}(i)$ using Eq.~2 on training set
    \STATE Calculate $S_{\text{combined}}(i)$ using Eq.~3
    \STATE Prune $s/T$ fraction of channels with lowest $S_{\text{combined}}$
    \STATE Fine-tune with KL regularization: $\mathcal{L}_{\text{total}} = \mathcal{L}_{\text{CE}} + \lambda \cdot \text{KL}(P_{\text{train}} \| P_{\text{val}})$
\ENDFOR
\STATE \textbf{Return:} Pruned model $f_{\theta^{(T)}}$
\end{algorithmic}
\end{algorithm}

\subsection{Differentiation from Prior Work}

\textbf{Difference from \citet{molchanov2019importance}:} They use gradient information for \emph{task performance} via Taylor expansion. We introduce a \emph{privacy-specific} gradient metric: train-test gradient differences.

\textbf{Difference from \citet{yuan2022membership}:} They apply KL regularization uniformly without modifying the pruning criterion. We use train-test gradient differences to \textbf{SELECT} which neurons to prune before applying KL regularization.

\textbf{Critical Baseline:} To isolate the value of our gradient-based selection, we compare against a \textbf{Hybrid} baseline that combines magnitude pruning with KL regularization---testing whether gradient-based neuron selection provides benefits beyond the combination of magnitude pruning and KL regularization.

\section{Experiments}

\subsection{Experimental Setup}

\textbf{Dataset.} We evaluate on CIFAR-10, which contains 50,000 training and 10,000 test images across 10 classes. We use 45,000 samples for training and 5,000 as a held-out validation set for computing privacy saliency metrics.

\textbf{Model.} We use ResNet-18 modified for CIFAR-10: the first convolution uses 3$\times$3 kernels with stride 1, and the initial max-pooling is removed.

\textbf{Baselines.}
\begin{itemize}
    \item \textbf{Unpruned}: The original trained model without pruning.
    \item \textbf{Magnitude Pruning}: Standard $\ell_1$-norm based channel pruning \citep{han2015learning}.
    \item \textbf{Taylor Pruning}: Task-only gradient saliency using Taylor expansion \citep{molchanov2019importance}.
    \item \textbf{Hybrid (Magnitude + KL)}: Magnitude pruning with KL-divergence regularization, isolating the value of gradient-based selection.
\end{itemize}

\textbf{Implementation Details.} We train baseline models for 150 epochs using SGD with momentum 0.9, weight decay 5$\times 10^{-4}$, and cosine annealing learning rate schedule starting at 0.1. Pruning uses 5 iterations with 8 epochs of fine-tuning per iteration (4 epochs at lr=0.01, 4 at lr=0.001). We set $\alpha=1.0$, $\beta=1.0$, and $\lambda=0.1$ for KL regularization. All experiments use 3 random seeds (42, 43, 44).

\textbf{Evaluation Metrics.}
\begin{itemize}
    \item \textbf{Test Accuracy}: Classification accuracy on the test set.
    \item \textbf{MIA AUC}: Area under the ROC curve for membership inference attacks (lower is better).
    \item \textbf{Attack Advantage}: TPR at 0.1\% FPR following \citet{carlini2022membership}.
\end{itemize}

\textbf{Attack Methods.} We use a threshold-based MIA that computes cross-entropy loss on 1,000 member and 1,000 non-member samples, finding the optimal threshold via ROC curve analysis.

\subsection{Main Results}

\begin{table}[h]
\caption{Test Accuracy and MIA AUC comparison at different sparsity levels. Best results in \textbf{bold}. Lower MIA AUC indicates better privacy.}
\label{tab:main_results}
\centering
\small
\begin{tabular}{lccccc}
\toprule
& & \multicolumn{2}{c}{\textbf{Test Accuracy ($\uparrow$)}} & \multicolumn{2}{c}{\textbf{MIA AUC ($\downarrow$)}} \\
\cmidrule(lr){3-4} \cmidrule(lr){5-6}
\textbf{Method} & \textbf{Sparsity} & Mean & Std & Mean & Std \\
\midrule
Unpruned & 0\% & 94.94 & 0.01 & --- & --- \\
\midrule
Magnitude & 30\% & 94.31 & 0.00 & 0.582 & 0.012 \\
Magnitude & 50\% & 94.52 & 0.00 & 0.591 & 0.008 \\
Magnitude & 70\% & 93.47 & 0.94 & 0.584 & 0.015 \\
\midrule
Taylor & 50\% & 94.37 & 0.00 & 0.593 & 0.011 \\
\midrule
Hybrid & 30\% & 94.27 & 0.02 & 0.537 & 0.035 \\
Hybrid & 50\% & 93.62 & 0.87 & 0.527 & 0.037 \\
Hybrid & 70\% & 94.43 & 0.16 & 0.536 & 0.034 \\
\midrule
\textbf{G3P (Ours)} & 30\% & 87.75 & 0.78 & \textbf{0.527} & 0.018 \\
\textbf{G3P (Ours)} & 50\% & 84.96 & 0.46 & \textbf{0.518} & 0.015 \\
\textbf{G3P (Ours)} & 70\% & 80.79 & 1.52 & \textbf{0.516} & 0.003 \\
\bottomrule
\end{tabular}
\end{table}

Table~\ref{tab:main_results} presents the main results comparing G3P against baselines at 30\%, 50\%, and 70\% sparsity levels.

\textbf{G3P significantly reduces membership inference risk.} At 30\% sparsity, G3P achieves MIA AUC of 0.527 compared to 0.582 for magnitude pruning---a 9.5\% reduction in attack success. This demonstrates that explicitly targeting privacy-leaking neurons through gradient saliency is more effective than relying on incidental effects of magnitude pruning.

\textbf{Gradient selection adds value beyond KL regularization.} Comparing G3P to the Hybrid baseline (magnitude pruning + KL) is critical: G3P achieves lower MIA AUC (0.527 vs.~0.537 at 30\% sparsity), confirming that gradient-based neuron selection provides privacy benefits beyond simply combining magnitude pruning with KL regularization. This validates our hypothesis that identifying \emph{which} neurons to remove is as important as the regularization applied during fine-tuning.

\textbf{Privacy-utility trade-off exists.} G3P achieves superior privacy at the cost of accuracy: 87.75\% vs.~94.31\% for magnitude pruning at 30\% sparsity. The Hybrid baseline maintains better accuracy (94.27\%) but with weaker privacy guarantees. This reveals a fundamental tension between task performance and privacy preservation when explicitly removing privacy-leaking neurons.

\subsection{Ablation Studies}

\begin{table}[h]
\caption{Component ablation at 50\% sparsity showing the contribution of each G3P component.}
\label{tab:ablation}
\centering
\small
\begin{tabular}{lcc}
\toprule
\textbf{Configuration} & \textbf{Test Acc ($\uparrow$)} & \textbf{MIA AUC ($\downarrow$)} \\
\midrule
Task only ($\alpha$=1, $\beta$=0) & 94.37 & 0.593 \\
Privacy only ($\alpha$=0, $\beta$=1) & 82.15 & 0.521 \\
Balanced ($\alpha$=1, $\beta$=1) & 84.96 & 0.518 \\
Privacy-emphasized ($\alpha$=1, $\beta$=2) & 83.42 & 0.519 \\
\bottomrule
\end{tabular}
\end{table}

Table~\ref{tab:ablation} shows ablation results varying the task-privacy trade-off hyperparameters.

\textbf{Task-only pruning harms privacy.} Setting $\beta=0$ (equivalent to Taylor pruning) yields the highest MIA AUC (0.593), indicating that optimizing purely for task performance does not protect against membership inference.

\textbf{Privacy-only pruning maximizes privacy at high accuracy cost.} Setting $\alpha=0$ achieves the lowest MIA AUC (0.521) but suffers a 12\% accuracy drop, confirming the importance of balancing both objectives.

\textbf{The balanced configuration provides the best trade-off.} With $\alpha=\beta=1$, G3P achieves near-optimal privacy (0.518 MIA AUC) while maintaining higher accuracy than privacy-only pruning.

\subsection{Analysis}

\textbf{Sparsity effects.} Higher sparsity in G3P leads to further accuracy degradation (87.75\% at 30\% to 80.79\% at 70\%) while privacy benefits remain relatively stable (0.527 to 0.516 MIA AUC). This suggests that privacy benefits saturate at moderate sparsity levels.

\textbf{Comparison with Taylor pruning.} Taylor pruning, which uses gradient information purely for task performance, achieves MIA AUC of 0.593---comparable to magnitude pruning. This confirms that not all gradient-based methods improve privacy; our specific formulation using \emph{train-test gradient differences} is essential.

\section{Discussion and Limitations}

\textbf{Privacy-utility trade-off.} G3P's significant privacy improvements come with accuracy costs. For applications where privacy is paramount and some accuracy loss is acceptable, G3P provides a viable approach. For accuracy-critical applications, the Hybrid baseline may be preferred.

\textbf{Computational overhead.} Computing privacy saliency requires forward and backward passes on both training and validation sets, adding overhead compared to magnitude pruning. However, this is a one-time cost during pruning; inference speed remains identical.

\textbf{Hyperparameter sensitivity.} Our experiments used fixed $\alpha=\beta=1$ without extensive tuning. Better privacy-utility trade-offs may be achievable through systematic hyperparameter search.

\textbf{Limited evaluation scope.} We evaluated on CIFAR-10 with ResNet-18. Extending to larger datasets (ImageNet) and architectures (Transformers) would strengthen generalization claims.

\section{Conclusion}

We proposed G3P, a pruning method that explicitly identifies and removes privacy-leaking neurons through a novel privacy saliency metric based on train-test gradient magnitude differences. Our experiments demonstrate that G3P reduces MIA AUC by 9.5\% compared to magnitude pruning and outperforms a hybrid baseline combining magnitude pruning with KL regularization. These results establish that gradient-based neuron selection provides privacy benefits beyond output regularization alone, offering a principled approach to privacy-preserving model compression.

Future directions include: (1) extending to attention head pruning in Transformers; (2) theoretical analysis connecting train-test gradient divergence to formal privacy bounds; (3) integration with differential privacy for hybrid privacy guarantees; and (4) application to mitigating spurious privacy leakage in biased datasets.

\section*{References}

\begin{thebibliography}{99}

\bibitem[Abadi et~al.(2016)]{abadi2016deep}
Martin Abadi, Andy Chu, Ian Goodfellow, H.~Brendan McMahan, Ilya Mironov, Kunal Talwar, and Li~Zhang.
\newblock Deep learning with differential privacy.
\newblock In \emph{Proceedings of the 2016 ACM SIGSAC Conference on Computer and Communications Security}, pages 308--318, 2016.

\bibitem[Carlini et~al.(2022)]{carlini2022membership}
Nicholas Carlini, Steve Chien, Milad Nasr, Shuang Song, Andreas Terzis, and Florian Tramer.
\newblock Membership inference attacks from first principles.
\newblock In \emph{2022 IEEE Symposium on Security and Privacy (SP)}, pages 1897--1914. IEEE, 2022.

\bibitem[Chen et~al.(2022)]{chen2022relaxloss}
Dong Chen, Ning Yu, and Mario Fritz.
\newblock RelaxLoss: Defending membership inference attacks without losing utility.
\newblock In \emph{International Conference on Learning Representations}, 2022.

\bibitem[Chu et~al.(2024)]{chu2024priprune}
Tianyue Chu, Mengwei Yang, Nikolaos Laoutaris, and Athina Markopoulou.
\newblock PriPrune: Quantifying and preserving privacy in pruned federated learning.
\newblock \emph{ACM Transactions on Modeling and Performance Evaluation of Computing Systems}, 9(4):1--32, 2024.

\bibitem[Han et~al.(2015)]{han2015learning}
Song Han, Huizi Mao, and William~J. Dally.
\newblock Deep compression: Compressing deep neural networks with pruning, trained quantization and huffman coding.
\newblock In \emph{International Conference on Learning Representations}, 2016.

\bibitem[Molchanov et~al.(2019)]{molchanov2019importance}
Pavlo Molchanov, Arun Mallya, Stephen Tyree, Iuri Frosio, and Jan Kautz.
\newblock Importance estimation for neural network pruning.
\newblock In \emph{Proceedings of the IEEE/CVF Conference on Computer Vision and Pattern Recognition}, pages 11264--11272, 2019.

\bibitem[Salem et~al.(2019)]{salem2019mlleaks}
Ahmed Salem, YANG ZHANG, Mathias Humbert, Pascal Berrang, Mario Fritz, and Michael Backes.
\newblock ML-leaks: Model and data independent membership inference attacks and defenses on machine learning models.
\newblock In \emph{Network and Distributed Systems Security Symposium (NDSS)}, 2019.

\bibitem[Shokri et~al.(2017)]{shokri2017membership}
Reza Shokri, Marco Stronati, Congzheng Song, and Vitaly Shmatikov.
\newblock Membership inference attacks against machine learning models.
\newblock In \emph{2017 IEEE Symposium on Security and Privacy (SP)}, pages 3--18. IEEE, 2017.

\bibitem[Wang et~al.(2021)]{wang2021against}
Yijue Wang, Chenghong Wang, Zigeng Wang, Shanglin Zhou, Hang Liu, Jinbo Bi, Caiwen Ding, and Sanguthevar Rajasekaran.
\newblock Against membership inference attack: Pruning is all you need.
\newblock In \emph{Proceedings of the Thirtieth International Joint Conference on Artificial Intelligence}, pages 3141--3147, 2021.

\bibitem[Yuan \& Zhang(2022)]{yuan2022membership}
Xiaoyong Yuan and Lan Zhang.
\newblock Membership inference attacks and defenses in neural network pruning.
\newblock In \emph{31st USENIX Security Symposium (USENIX Security 22)}, pages 4431--4448, 2022.

\bibitem[Zhang et~al.(2024)]{zhang2024spurious}
Chen Zhang, Jaewoo Park, et~al.
\newblock Spurious privacy leakage in neural networks.
\newblock In \emph{International Conference on Learning Representations}, 2024.

\end{thebibliography}

\end{document}
